\section{Comparison data, the toric mechanism and the Hamiltonian}\label{p2:sec-machine}

This section states the comparison data for concave profiles and the packing theorem in a toric model, both proved
in the supplement, and recalls from Paper~I \cite{PaperI} the toric balls, the twist and the passage from embeddings
to one Hamiltonian. Statements recalled from Paper~I carry the tag ``Paper~I'' and have the same name there. Here
$m=n-1$; the comparison data serve the higher-dimensional models and, for $m=1$, the scalar branch in dimension four.

\emph{Dependencies.} Theorem~\ref{thm:cmp-main} follows from Proposition~\ref{prop:cmp-fixed-point} and
Theorem~\ref{thm:cmp-realisation}; the proofs are in \Ssec{sec:fcp} and \Ssec{sec:plc}. Theorem~\ref{thm:mech-packing}
is \Sref{Theorem}{thm:shd-packing}, with the handle construction of \Ssec{p2:sec-shd}. The termination part of
Theorem~\ref{thm:eff-algorithm} uses Theorem~\ref{thm:main}(i), proved in \S\ref{p2:sec-proof}. The statements of
the supplement of Paper~I used in the proofs are reproduced in \Ssec{p2:sec-quoted}.

\emph{Conventions.} The dilation $z\mapsto\lambda z$ multiplies $\omega_0$ and capacities by $\lambda^2$; conjugation by
it normalises the target capacity to $1$. In fibre coordinates $w_j$ the moments are $p_j=\pi|w_j|^2$ and the angles
$\theta_j=\arg w_j/2\pi\in\R/\Z$, so that the fibre form is $\omega_V=\sum_jdp_j\wedge d\theta_j$ away from the axes.
Fibre Hamiltonians $Q$ follow the sign of \cite{OpenAI2026}, $\iota_{X_Q}\omega_V=-dQ$, opposite to the global one. In
affine coordinates $\omega_{\rm FS}=dd^c\log(1+|z|^2)$ with $d^c=i(\bar\partial-\partial)/(4\pi)$; the map
$z\mapsto z/\sqrt{\pi(1+|z|^2)}$ identifies this chart symplectically with $\Bo(1)$. All cutoffs and clocks use the flat
step $S$ of \Sref{Definition}{I:def:pre-flat}, all of whose derivatives vanish at $0$ and $1$.

\subsection{Toric balls and the twist}\label{I:sec:mech-toric}

A downward polytope is
\[
 \Delta=\{p\in\R^m_{\ge0}:\langle b_\nu,p\rangle\le\ell_\nu,\ 1\le\nu\le N\},
 \qquad b_\nu\in\R^m_{\ge0},\quad \ell_\nu>0,
\]
with toric domain $V_\Delta=\{w\in\C^m:p(w)\in\Delta\}$. The equalities $\langle b_\nu,p\rangle=\ell_\nu$ are its outer
faces, and $\tau_\Delta\Delta$, $0<\tau_\Delta<1$, is a core that keeps the source caps away from them.

\begin{definition}[Toric ball; \PoneTag{def:mech-toric-ball}]\label{I:def:mech-toric-ball}
\lean{OAI.PackingSufficiencySupport.closedBall,
OAI.PackingSufficiencySupport.Hamiltonian.splitCapacityBall,
OAI.PackingSufficiencySupport.Hamiltonian.complexSplitPhase_capacity}
In $\C^n=\C\times\C^m$, put $\varpi=\pi|z_{\rm b}|^2$. The ball of capacity $a$ is
$B(a)=\{\varpi+|p|\le a\}$. At a fibre point with $|p|\le a$ its base slice is a disc of capacity $a-|p|$, so its
slice-height function is $(a-|p|)_+$.
\end{definition}

\begin{lemma}[The cap integral; \PoneTag{lem:mech-toric-volume}]\label{I:lem:mech-toric-volume}
\lean{OAI.PackingSufficiencySupport.integral_simplex_tent, OAI.PackingSufficiencySupport.integral_tent}
For $a\ge0$, $\int_{\R^m_{\ge0}}(a-|p|)_+\,dp=a^{m+1}/(m+1)!$.
\end{lemma}
\begin{proof}
The simplex $\{|p|\le t\}$ has measure $t^m/m!$; integrate its distribution function against $(a-t)_+$.
\end{proof}

Each model has a compact affine core $\Sigma_\ell$ and a horizontal primitive $\Gamma_{\rm phys}(p)$; the physical form is
$\Omega=\omega_V+d\Gamma_{\rm phys}$, with base curvature $\kappa(p)=d_\Sigma\Gamma_{\rm phys}(p)$. The base profile and budget
are
\begin{equation}\label{I:eq:mech-model-volume}
 \Hb(p)=\int_{\Sigma_\ell}\kappa(p),\qquad
 \mathcal B_\Delta=\int_\Delta\Hb(p)\,dp .
\end{equation}
For auxiliary capacities $a_i$ larger than the given ones, with $\{|p|\le a_i\}\subset\tau_\Delta\Delta$, put
$\Pb(p)=\Hb(p)-\sum_i(a_i-|p|)_+$. By the cap integral,
\begin{equation}\label{I:eq:mech-mean}
 \int_\Delta\Pb=\mathcal B_\Delta-\sum_i\frac{a_i^n}{n!},\qquad
 \frac{1-\sum_i\rho_i^n}{n!}
 =\Big(\frac1{n!}-\mathcal B_\Delta\Big)+\int_\Delta\Pb+\sum_i\frac{a_i^n-\rho_i^n}{n!}.
\end{equation}
For small balls the model parameters make the first and last terms on the right smaller than the left side, so
$\int_\Delta\Pb>0$; for one large ball of reserved capacity $a$ the limiting budget for the remaining balls is
$(1-a^n)/n!$. A positive integral does not give the pointwise inequality $\Pb\ge0$ needed to stack the base discs; the
twist uses only the integral.

\emph{The twist.} The handle data (\Sref{Definition}{def:hdl-handle-data}) include an annulus with coordinates
$(s,v)\in[0,1]\times\R/\Z$, a closed one-form $\beta=b(v)\,dv$ with $b\ge0$, $\int_0^1b=1$, and a positive reference
area form $\Lambda$. A first Moser path replaces the physical curvature on the annulus by a background
$e_\Lambda\Lambda$ plus the base mass $H^*(p)\ge\Hb(p)-\gamma_*/16$ carried by $\beta$
(\Ssec{sec:hdl-first}), where $\gamma_*$ is a reserve set aside from the gap below. Let $q_t$ be a Hamiltonian
isotopy of the fibre, $q=q_1$; the suspension $\mathcal D_\lambda(s,v,w)=(s,v,q_{\lambda f(s)}w)$, $f=S$, inserts it
across the annulus. For a function $K$ of the moments put $J=K-H^*$; adding $d(J(p)\beta)$ changes only toric mixed
terms away from the annulus, and after the twist the base area available across the annulus at $w$ is
\begin{equation}\label{I:eq:mech-available}
 K\big(p(q^{-1}w)\big)-J\big(p(w)\big)
 =H^*\big(p(w)\big)+K\big(p(q^{-1}w)\big)-K\big(p(w)\big).
\end{equation}
Since $q$ preserves the fibre volume, this redistributes area between moment levels. A comparison supplies, for a
minorant $P\le\Pb$,
\begin{equation}\label{I:eq:mech-comparison}
 (K-P)\big(p(qw)\big)\le K\big(p(w)\big)-\gamma\qquad(w\in V_\Delta),
 \qquad 0<\gamma_*\le\gamma.
\end{equation}
Applied at $q^{-1}w$ it shows that the available area exceeds $\sum_i(a_i-|p(w)|)_+$ by at least
$\gamma-\gamma_*/16>0$. This is divided among disjoint \emph{layers} $I_1,\dots,I_k,I_*$ of $s\in(0,1)$: across $I_i$
the $\beta$-coefficient rises by a smooth cap $h_i(p)$ below $(a_i-|p|)_+$, and a ball of capacity slightly less than
$a_i$ is placed there in product form; across $I_*$ it rises by the positive remainder. The general comparator is
built from a concave piecewise affine minorant $P$ with positive mean, with
\begin{equation}\label{I:eq:mech-necessary}
 0<d<\min\bigl(\bar P,\beta_P\bigr),\qquad
 \bar P=|\Delta|^{-1}\int_\Delta P,
\end{equation}
where $P\ge\beta_P>0$ off the core $\tau_\Delta\Delta$ (\S\ref{p2:sec-machine-cmp}).

\subsection{Model, comparison datum and transfer}\label{I:sec:mech-interfaces}

\begin{definition}[Model input]\label{def:mech-model}
\lean{SymplecticPacking.HighDim.RawToricModel,
SymplecticPacking.HighDim.ToricPackingStages.smallRawModel,
SymplecticPacking.HighDim.ToricPackingStages.mixedRawModel,
SymplecticPacking.HighDim.ToricPackingStages.cubicRawTotalModel}
A higher-dimensional model input consists of a compact affine core $\Sigma_\ell\subset M$ of an affine
curve $M$, the product $E=M\times\C^m$, a downward polytope $\Delta$, and the physical primitive
$\Gamma_{\rm phys}(p)$ given in the corresponding model section, together with the handle data of
\Sref{Definition}{def:hdl-handle-data}. The physical form is
\[
 \Omega=\omega_V+d\Gamma_{\rm phys},\qquad
 V_\Delta(E)=\Sigma_\ell\times V_\Delta.
\]
The physical curvature $\kappa(p)=d_\Sigma\Gamma_{\rm phys}(p)$ is positive on the model region and extends
smoothly in the moment variables across the coordinate hyperplanes.
\end{definition}

\begin{definition}[Comparison datum]\label{def:mech-datum}
\lean{SymplecticPacking.HighDim.ComparisonDatum,
SymplecticPacking.HighDim.ComparisonDatum.qMap,
SymplecticPacking.HighDim.CmpInput.comparatorDatum}
A comparison datum for $(\Delta,P)$ with gap $\gamma>0$ and reserve $\gamma_*\in(0,\gamma]$ consists of a function
$K$ of the moments, smooth on a neighbourhood of $\Delta$ in $\R^m$, and a smooth Hamiltonian isotopy
$(q_t)_{t\in[0,1]}$ of $\C^m$ with $q_0=\mathrm{id}$, generated by $Q_t$ with $\iota_{X_{Q_t}}\omega_V=-dQ_t$,
such that:
\begin{enumerate}[label=(\arabic*),ref=(\arabic*)]
\item\label{it:mech-datum-collar} $q_t(V_\Delta)=V_\Delta$ for every $t$, and there is $\varsigma>0$ such that
 for every face $\nu$ and every $t$, $Q_t$ Poisson-commutes with the face moment $\mu_\nu=\langle b_\nu,p\rangle$ on
 the collar $\{\mu_\nu>\ell_\nu-\varsigma\}$;
\item\label{it:mech-datum-cmp} $q=q_1$ satisfies the comparison inequality~\eqref{I:eq:mech-comparison}.
\end{enumerate}
For the comparator of Theorem~\ref{thm:cmp-main}, with gap $\gamma=\frac{11}{16}d$, the models take the reserve
$\gamma_*=\frac{11}{64}d$.
\end{definition}

\begin{definition}[Transfer]\label{def:mech-transfer}
\lean{SymplecticPacking.BranchSmall.vectorTransfer,
SymplecticPacking.BranchLargeHigh.Two.returnMap,
SymplecticPacking.BranchLarge3.RawTransfer.transferMap,
SymplecticPacking.Return.compositeT}
A transfer to an open target $U\subset\C^n$ assigns to a compact subset $\mathcal K$ of the model
domain a smooth symplectic embedding $\mathcal T_{\mathcal K}$ on a neighbourhood of $\mathcal K$,
and an explicit compact set $\mathcal C_{\mathcal K}\subset U$ containing the image of $\mathcal K$.
The cubic, quadric and six-dimensional polynomial transfers are fixed maps; in the shell transfer for
$n\ge4$ the slit-shell parameters depend on $\mathcal K$.
\end{definition}

\begin{theorem}[Packing in a model]\label{thm:mech-packing}
\lean{SymplecticPacking.HighDim.RawModelAgreement.native_packing_spec, SymplecticPacking.Strip.Cubic.fN_symplectic, SymplecticPacking.Strip.Cubic.fN_disjoint, SymplecticPacking.Strip.Cubic.fN_avoid_fibres}
Let $n\ge3$ and let $(\Sigma_\ell,\Delta,\Gamma_{\rm phys})$ be a model input
(Definition~\ref{def:mech-model}) whose handle data satisfy the collar bounds (\Snum{eq:hdl-collar}). Let
$(K,(q_t))$ be a comparison datum for $(\Delta,P)$ with gap $\gamma$ and reserve $\gamma_*$
(Definition~\ref{def:mech-datum}). Let $0<\rho_i<a_i$ $(1\le i\le k)$ with
$\{|p|\le a_i\}\subset\tau_\Delta\Delta$, and suppose
\[
 P(p)\le\Hb(p)-\sum_i(a_i-|p|)_+\qquad(p\in\Delta).
\]
Then the maps $\psi_i$ of \Sref{Construction}{cons:native-packing} are smooth symplectic embeddings of
neighbourhoods of the closed balls $B(\rho_i)$ into $(E,\Omega)$, and the sets $\psi_i(B(\rho_i))$ are pairwise
disjoint and contained in $\operatorname{int}V_\Delta(E)$. For $n=2$ the same conclusion holds for the cubic model of
Paper~I with its handle data and equivariant maps, when the comparison datum is replaced by valid cubic data
(\PoneRef{def:shd4-data}): the full-disc comparison for the profile itself and the collar condition $P_{\rm supp}<h$.
\end{theorem}

The maps are $\psi_i=\Phi^{-1}\circ\mathcal D^{-1}\circ\mathcal E_i$: a layer map $\mathcal E_i$ placing the ball in the
layer $I_i$, the inverse twist, and the inverse of three Moser maps (\Ssec{sec:hdl-proof}). The proof for the
bordered base is \Sref{Theorem}{thm:shd-packing}, with \Sref{Propositions}{prop:str-layer}
and~\Snum{prop:str-embeddings}. The model sections verify the density floor, the collar bounds and the profile
minorant before applying the theorem.

\begin{proposition}[Assembly]\label{prop:mech-assembly}
\lean{OAI.PackingSufficiencySupport.Hamiltonian.postcompose_form_packing, SymplecticPacking.HighDim.assembly_spec}
Let $\psi_i$ be the model embeddings of Theorem~\ref{thm:mech-packing} and let $\mathcal T$ be its transfer.
Then, with $\mathcal K=\bigcup_i\psi_i(B(\rho_i))$, the maps
$f_i=\mathcal T_{\mathcal K}\circ\psi_i$ are smooth symplectic embeddings on neighbourhoods of
$B(\rho_i)$, have pairwise disjoint images in $U$, and have images contained in $\mathcal C_{\mathcal K}$.
\end{proposition}
\begin{proof}
$\mathcal K$ is compact and the transfer is one injective symplectic map on a neighbourhood of it, so composition
preserves the pullback identity and disjointness; the inverse images of that neighbourhood under the $\psi_i$ are open
neighbourhoods of the source balls.
\end{proof}

\begin{table}[ht]
\centering\footnotesize
\renewcommand{\arraystretch}{1.2}
\begin{tabular}{@{}L{.16\textwidth}L{.24\textwidth}L{.25\textwidth}L{.25\textwidth}@{}}
\toprule
input & base and physical primitive & base profile & target\\\midrule
$n=2$ (Paper~I) & disc bundle over the cubic, on its periodic cover & $3(1-3p)$ & $\Bo(1)$\\
$n\ge3$, all $c_i<C/2$ & $y_j^2=t^2-j$, $m=n-1$; $(1-2|p|)\lambda$ & $A_\ell(1-2|p|)$; limit $2^m(1-2|p|)$ & $\Bo(1)$\\
$n\ge4$, one large ball & the curve $y_j^2=t^2-j$, $j\le n-2$, fibre $\C^{n-1}$; $(1-u-2w)\lambda$ & $A_\ell(1-u-2w)$; limit $2^{n-2}(1-u-2w)$ & a shell about the reserved ball\\
$n=3$, one large ball & an explicit polynomial primitive on a compact affine core & $H_\ell$, limit $9(b-|p|)-2(s-|p|)_+$ & common polynomial form, then $\Bo(1)$\\\bottomrule
\end{tabular}
\caption{The models. $A_\ell$ is the reference area of the compact affine core; each
higher-dimensional total space is a product over that core. For one large ball, $u=p_1$ and $w=p_2+\dots+p_{n-1}$
($n\ge4$); for $n=3$, $H_\ell$ is the base profile of the truncated polynomial model, with limit $H_0$, and
$s<b$ are explicit parameters of that model (\Ssec{sec:lrg-six}).}
\label{tab:mech-models}
\end{table}

\subsection{Comparison data}\label{p2:sec-machine-cmp}

Let $\Delta$ be a compact full-dimensional downward polytope and $P$ a continuous concave function on $\Delta$ with
positive mean. Norms are uniform norms on $\Delta$. For a coordinate $j$ write $p=(y,x)$ with $x=p_j$; then
$\Delta=\{(y,x):y\in\Delta_j,\ 0\le x\le L(y)\}$ with
$L(y)=\min_{\nu:\,b_{\nu j}>0}(\ell_\nu-\sum_{i\ne j}b_{\nu i}y_i)/b_{\nu j}$ concave piecewise affine on the projection
$\Delta_j$. The polytope is \emph{rational} if all $b_\nu,\ell_\nu$ are rational; a function on $\Delta$ is
\emph{rational convex polyhedral} (resp.\ \emph{rational concave piecewise affine}) if it is the maximum (resp.\
minimum) of finitely many affine functions with rational coefficients.

\begin{definition}[Coordinate rearrangement]\label{def:cmp-rearrangement}
For a continuous convex function $F$ on $\Delta$ and a coordinate $j$, put
\[
 (R_jF)(y,x)=\min_{\substack{0\le a,\ a+x\le L(y)}}\ \max\bigl(F(y,a),\,F(y,a+x)\bigr),
 \qquad (y,x)\in\Delta,
\]
and $\Rcal=R_m\circ\dots\circ R_1$.
\end{definition}

Thus $x\mapsto(R_jF)(y,x)$ is the increasing rearrangement of $F(y,\cdot)$ on $[0,L(y)]$; its properties
(monotonicity, nonexpansiveness, and preservation of rational convex polyhedral functions by Fourier--Motzkin
elimination) are \Sref{Lemma}{lem:cmp-rearrangement}.

\begin{proposition}[Finite comparator]\label{prop:cmp-fixed-point}
Let $\Delta$ be rational, $\tau_\Delta\in(0,1)$ rational, and $P$ a rational concave piecewise affine function on
$\Delta$ with mean $\mu=\vol(\Delta)^{-1}\int_\Delta P>0$ and $P\ge\beta_P$ on $\Delta\setminus\tau_\Delta\Delta$ for some
$\beta_P>0$. Let $d$ be rational with $0<d<\min(\mu,\beta_P)$ and put
\[
 \eta_\Delta=(1-\tau_\Delta)(\mu-d),\qquad M_P=\max\bigl(1,\ d-\min_\Delta P+1\bigr),\qquad
 \lambda=1-\frac{\eta_\Delta}{2(M_P+\eta_\Delta)}.
\]
Define $K_0=0$, $K_{r+1}=\Rcal(\lambda K_r-P)$ and
\[
 N=\left\lceil\frac{16\|P\|}{d(1-\lambda)(\lambda^{-1}-1)}\right\rceil+1.
\]
Then every $K_r$ is rational convex polyhedral, $\lambda^N\|P\|/(1-\lambda)<d/16$, and
\[
 \Rcal(K_N-P)\le K_N-\tfrac78d\qquad\text{on }\Delta .
\]
\end{proposition}

All quantities are computed by finite rational arithmetic and Fourier--Motzkin elimination. The construction, and the
smoothing and realisation below, apply verbatim over real coefficients (``convex polyhedral'' in place of ``rational
convex polyhedral''); this version is used for real capacities.

\begin{theorem}[Realisation of one coordinate rearrangement]\label{thm:cmp-realisation}
\lean[needed]{SymplecticPacking.HighDim.liftGen_spec}
\uses{lem:cmp-rearrangement, S-def:plc-slice, S-thm:plc-lift, S-thm:pln-isotopy}
Let $\Delta$ be rational, $F$ rational convex polyhedral on $\Delta$, $j\in\{1,\dots,m\}$ and $e>0$ rational. Let
$Q^{(j)}_t$ be the Hamiltonian defined from $(\Delta,F,j,e)$ in \Sref{Definition}{def:plc-slice}, and let
$(q^{(j)}_t)_{t\in[0,1]}$ be its flow ($\iota_{X}\omega_V=-dQ^{(j)}_t$). Then $Q^{(j)}$ is smooth with compact support
in $w$, $(q^{(j)}_t)$ is a smooth Hamiltonian isotopy of $\C^m$, and:
\begin{enumerate}[label=(\alph*)]
\item $q^{(j)}_0=\mathrm{id}$; every $q^{(j)}_t$ maps $V_\Delta$ onto itself and preserves the moments $p_i$, $i\ne j$.
\item There is a rational $\varsigma_j>0$ such that for every $\nu$ with $b_{\nu j}>0$ and every $t$, $Q^{(j)}_t$
 vanishes on $\{\langle b_\nu,p\rangle>\ell_\nu-\varsigma_j\}$; for every $\nu$ with $b_{\nu j}=0$, $Q^{(j)}_t$
 Poisson-commutes with $\langle b_\nu,p\rangle$.
\item $F\bigl(p(q^{(j)}_1(w))\bigr)\le(R_jF)(p(w))+e$ for every $w\in V_\Delta$.
\end{enumerate}
\end{theorem}

The proof is \Sref{Theorem}{thm:plc-lift}. The construction is planar: fix $y$ and rescale an inner slice
$[0,\ell_{\rm in}(y)]$, $\ell_{\rm in}(y)<L(y)$, to $[0,1]$; then $f(v)=F(y,\ell_{\rm in}(y)v)$ is convex with increasing
rearrangement $f^\sharp$, and one needs an area-preserving map $q_{\rm in}$ of the disc with
$f(\pi|q_{\rm in}(z)|^2)\le f^\sharp(\pi|z|^2)+\xi$ on a core disc. The sublevel sets of a convex potential approximating
$f(v)$ (log-sum-exp smoothing plus a barrier), proper in the angle and exactly quadratic near its minimum, give such a
map through Morse charts (\Sref{Theorem}{thm:plc-planar}). A compactly supported Hamiltonian isotopy of the disc ends
at this map on a core (\Sref{Theorem}{I:thm:pln-isotopy}); its constants are uniform in $y$
(\Sref{Theorem}{thm:plc-param}), and the planar generator is lifted as
\[
 Q^{(j)}_t(w)=\chi(y)S'(t)\ell_{\rm in}(y)\,
 Q^{\rm pl}\bigl(y,\chi(y)S(t),w_j/\sqrt{\ell_{\rm in}(y)}\bigr),
\]
with a cutoff $\chi(y)$ vanishing near degenerate slices.

\begin{theorem}[Comparison datum for concave profiles]\label{thm:cmp-main}
\uses{prop:cmp-fixed-point, thm:cmp-realisation, def:mech-datum, S-lem:fcp-smoothing}
Assume the hypotheses of Proposition~\ref{prop:cmp-fixed-point}, and let $d$, $N$ and $K_N$ be as there. Let
$K_N=\max_{\iota\in I}\phi_\iota$ be the representation by rational affine functions produced by the iteration, each
step $R_j$ being the Fourier--Motzkin elimination of \Sref{Lemma}{lem:cmp-rearrangement}(v), set
$c_K=\lceil32\,\#I/d\rceil+1$ (so that $\log(\#I)/c_K<d/32$ by \Sref{Lemma}{lem:fcp-smoothing}), and put
\[
 K=\frac1{c_K}\log\sum_{\iota\in I}e^{c_K\phi_\iota}.
\]
Let $F_0=K_N-P$ and $F_j=R_jF_{j-1}$ ($1\le j\le m$); let $q^{(j)}_t$ be the isotopy of
Theorem~\ref{thm:cmp-realisation} for $F_{j-1}$, the coordinate $j$ and $e_j=d/(8(m+1))$. Define $q_t$ by running
$q^{(m)},q^{(m-1)},\dots,q^{(1)}$ in turn with flat clocks: for $t\in[\frac{r-1}m,\frac rm]$, $r=1,\dots,m$, and
$j=m+1-r$,
\[
 q_t=q^{(j)}_{S(mt-r+1)}\circ q^{(j+1)}_1\circ\dots\circ q^{(m)}_1 .
\]
Then $(K,(q_t))$ is a comparison datum for $(\Delta,P)$ with gap $\gamma=\tfrac{11}{16}d$.
\end{theorem}

The proof chains the realisation errors $e_j$ along $u_m=w$, $u_{j-1}=q^{(j)}_1(u_j)$, which gives
$(K_N-P)(p(q_1(w)))\le K_N(p(w))-\frac34d$, and uses $K_N\le K<K_N+d/32$.

\subsection{From embeddings to one Hamiltonian}\label{p2:sec-machine-ham}

The models give explicit symplectic embeddings of slightly enlarged balls centred at the origin into $\Bo(1)$, with
pairwise disjoint images. Paper~I joins the translations $u\mapsto s_i+u$ to such embeddings by one explicit family of
symplectic embeddings with pairwise disjoint images at every time, and writes down a Hamiltonian generating it on the
moving balls. Both steps hold in every dimension $n\ge1$.

\begin{definition}[Normalisation; \PoneTag{def:ham-normalise}]\label{I:def:ham-normalise}
Let $C>0$ and $c\in(0,\infty)^k$. The \emph{normalised capacities} are $\rho_i=c_i/C$. For a map $f^\natural$
defined on a subset of $\R^{2n}$, its \emph{physical rescaling} is $f(u)=\sqrt C\,f^\natural(u/\sqrt C)$.
\end{definition}

Admissibility is homogeneous, so $(n,k,C,c)$ is admissible if and only if $(n,k,1,\rho)$ is; the physical rescaling
is symplectic, injective or of disjoint images exactly when $f^\natural$ is.

\begin{definition}[Embedding datum; \PoneTag{def:ham-datum}]\label{I:def:ham-datum}
\lean{SymplecticPacking.General.EmbeddingData, SymplecticPacking.General.EmbeddingSpec}
Let $n\ge1$, $k\ge1$, $C>0$ and $c\in(0,\infty)^k$. An \emph{embedding datum} for $(n,k,C,c)$ consists of numbers
$c'_i,\bar c_i$ with $c_i<c'_i<\bar c_i$ and maps $f_i\colon\Bo(\bar c_i)\to\R^{2n}$ such that
\begin{enumerate}[label=(\roman*)]
\item $f_i$ is $C^\infty$ and $Df_i(u)\in\operatorname{Sp}(2n)$ for all $u\in\Bo(\bar c_i)$;
\item $f_i$ is injective on $B(c'_i)$;
\item the sets $f_i(B(c'_i))$, $1\le i\le k$, are pairwise disjoint and contained in $\Bo(C)$.
\end{enumerate}
\end{definition}

Given an embedding datum and centres $s_i$ with $B(s_i,c_i)$ pairwise disjoint, Paper~I fixes radii
$r_i=\sqrt{c_i/\pi}<r_{1,i}<r_{2,i}<r_{3,i}<r_{4,i}<r'_i=\sqrt{c'_i/\pi}$ (\PoneRef{lem:ham-radii}) and defines an
isotopy $F_{i,t}$ by explicit stages: a collision-free route of translations, a path in $\operatorname{Sp}(2n)$ from $I$ to
$Df_i(0)$, a shape stage, and the Alexander tail $u\mapsto v^{-1}f_i(vu)$ (\PoneRef{cons:iso-isotopy}). Here
$\EB(y,r)$ and $\EB^\circ(y,r)$ are the closed and open Euclidean balls, and $\EB(r)=\EB(0,r)$.

\begin{proposition}[The common symplectic isotopy; \PoneTag{prop:ham-isotopy}]\label{I:prop:ham-isotopy}
\lean{SymplecticPacking.General.Ham.Isotopy.IsoInput.commonIsotopy, SymplecticPacking.General.Ham.Isotopy.IsoInput.contDiffOn_isoF, SymplecticPacking.General.Ham.Isotopy.IsoInput.isoF_sympl, SymplecticPacking.General.Ham.Isotopy.IsoInput.isoF_injOn, SymplecticPacking.General.Ham.Isotopy.IsoInput.isoF_disj, SymplecticPacking.General.Ham.Isotopy.IsoInput.isoF_bounded, SymplecticPacking.General.Ham.Isotopy.TargetIsoInput.isoF_norm_le}
The isotopy $F_{i,t}$ has the following properties.
\begin{enumerate}[label=(\alph*)]
\item $(t,u)\mapsto F_{i,t}(u)$ is $C^\infty$ on $\R\times\EB^\circ(r_{4,i})$; $F_{i,0}(u)=s_i+u$, $F_{i,1}=f_i$;
 $F_{i,t}=F_{i,0}$ for $t\le0$ and $F_{i,t}=F_{i,1}$ for $t\ge1$;
\item $D_uF_{i,t}(u)\in\operatorname{Sp}(2n)$ for all $t$ and $u\in\EB^\circ(r_{4,i})$;
\item each $F_{i,t}$ is injective on $\EB(r_{3,i})$, and for every $t$ the sets $F_{i,t}(\EB(r_{3,i}))$, $1\le i\le k$, are
 pairwise disjoint;
\item $F_{i,t}(\EB(r_{3,i}))\subset\EB(R_F)$ for all $i,t$, for an explicit radius $R_F$.
\end{enumerate}
\end{proposition}

\begin{definition}[The Hamiltonian of an embedding datum; \PoneTag{def:ham-H}]\label{I:def:ham-H}
\lean{SymplecticPacking.General.Ham.radialCutoff, SymplecticPacking.General.Ham.pushforward, SymplecticPacking.General.Ham.ofIsotopy, SymplecticPacking.HamEnd.hamiltonianOfEmbeddings}
For $u\in\EB^\circ(r_{4,i})$ let $\mathfrak a_{i,t}(u)[w]=\omega_0\big(\dot F_{i,t}(u),D_uF_{i,t}(u)w\big)$ and
$\mathfrak h_{i,t}(u)=\int_0^1\mathfrak a_{i,t}(\vartheta u)[u]\,d\vartheta$. Let $\chi_i=\chi_{r_{1,i},r_{2,i}}$ be the radial
cutoff of \Sref{Definition}{I:def:pre-cutoff}. Put $H_i(t,y)=\chi_i(u)\,\mathfrak h_{i,t}(u)$ if $u$ is the unique point of
$\EB(r_{3,i})$ with $F_{i,t}(u)=y$, and $H_i(t,y)=0$ if there is no such unique point; and $H_{\mathfrak E}=\sum_{i=1}^kH_i$,
where $\mathfrak E$ denotes the embedding datum together with the centres.
\end{definition}

\begin{theorem}[The Hamiltonian realises the isotopy; \PoneTag{thm:ham-flow}]\label{I:thm:ham-flow}
\lean{SymplecticPacking.General.packingFlow_diffeomorphism, SymplecticPacking.General.packingFlow_integral}
Let $H=H_{\mathfrak E}$. Then $H$ is $C^\infty$, vanishes for $t\notin(0,1)$ and outside $\R\times\EB(R_F)$; the flow
$\phi^t_H$ is defined for all $t\in\R$ and consists of symplectic diffeomorphisms of $\R^{2n}$, and
\[
 \phi^t_H(s_i+u)=F_{i,t}(u)\qquad(u\in\EB(r_i),\ t\in\R,\ 1\le i\le k).
\]
In particular $\phi^1_H(B(s_i,c_i))=f_i(B(c_i))$, and these sets are pairwise disjoint and contained in $\Bo(C)$.
\end{theorem}

This combines \PoneRef{lem:ham-smooth} with \PoneRef{thm:ham-flow}: the one-form $\mathfrak a_{i,t}$ is closed
because $F_{i,t}$ is symplectic, $\mathfrak h_{i,t}$ is its radial primitive, and on $F_{i,t}(\EB^\circ(r_{1,i}))$ the field
$X_H$ equals $\dot F_{i,t}$.

\begin{proposition}[Hamiltonian of an embedding datum; \PoneTag{prop:ham-end}]\label{I:prop:ham-end}
\lean{SymplecticPacking.HamEnd.hamiltonian_end}
Let $n\ge1$, let $(c'_i,\bar c_i,f_i)$ be an embedding datum for $(n,k,C,c)$, and let $s_1,\dots,s_k$ be centres such
that the closed balls $B(s_i,c_i)$ are pairwise disjoint. Then the Hamiltonian $H_{\mathfrak E}$ of
Definition~\ref{I:def:ham-H} satisfies $\Pack(n,k,C,c,s)$.
\end{proposition}

\subsection{Explicitness}\label{p2:sec-machine-eff}

\begin{definition}[Levels of explicitness]\label{def:intro-levels}
\begin{enumerate}[label=(L\arabic*)]
\item An \emph{explicit definitional construction} is a definition built from rational constants and the
 input by arithmetic and elementary functions, integrals over compact domains, limits of series with an
 explicit rate of convergence, maxima and minima of continuous functions over explicit compact sets,
 unique solutions of explicit ordinary differential equations, inverses of explicit injective maps,
 unique roots of explicit equations, case distinctions on propositions about the input and about the arguments of
 the function being defined, and least
 elements of searches over a fixed enumeration of finitely described objects (integers, rationals, and
 tuples of these and of rational polynomials) whose validity is decided by finitely many explicit
 inequalities; definite descriptions are used only after unique existence has been proved
 (\PoneRef{def:intro-the}).
\item The \emph{finite form} is $\Hfin=\beta_B\cdot P$ with $B\in\N$ and $P\in\Q[t,z]$, where
 $\beta_B(z)=\prod_{l=1}^{2n}\beta(\zeta_l/B)$ for the real coordinates $\zeta_1,\dots,\zeta_{2n}$ of $z$,
 and $\beta(x)=1-S\big((x^2-\frac14)/\frac5{16}\big)$; thus $\beta\colon\R\to[0,1]$ is smooth, equal to $1$ on
 $[-\frac12,\frac12]$ and to $0$ outside $(-\frac34,\frac34)$.
\item \emph{Computability} means that one oracle algorithm, given names of the input (a rational input is named by
 its constant sequence), halts on every admissible input and outputs some $(B,P)$ as in (L2) for which $\Hfin$
 satisfies the packing specification. The output may depend on the names and not only on the numbers they name;
 this is necessary, as explained below.
\end{enumerate}
\end{definition}

The algorithm of Section~\ref{sec:eff} accepts the first rational polynomial whose Euler trajectories from a finite grid,
with an error bound valid on $[0,1]$, stay where the cutoff equals one and end inside the target; every accepted output
is correct by a soundness theorem. Termination needs only some Hamiltonian with a positive margin; the existence
theorems would also supply one, non-constructively. We use $H[n,k,C,c,s]$.

The dependence on names in L3 is necessary for finite output. A computable single-valued map to the discrete set
of polynomial codes is continuous, hence constant on $\eps\in(0,1)$ for fixed $k$. But the images of $k\ge2$
disjoint closed balls under any fixed Hamiltonian diffeomorphism form a disconnected compact set, so their density in
their smallest enclosing ball is strictly below one, and that fixed output fails the requirement for sufficiently
small $\eps$.

